\clearpage

\begin{figure}[H]
\centering
\fbox{%
\begin{minipage}{0.92\columnwidth}
\small
\textbf{Original prompt}

\medskip
\textbf{System:} You are an evaluator of expert-domain scientific writing.
You will receive a query--answer pair, evaluation criteria, a scoring rubric,
and examples demonstrating how the evaluation should be performed. First
output your reasoning between \texttt{<reasoning>} and
\texttt{</reasoning>}, followed by the numeric score between
\texttt{<score>} and \texttt{</score>}.

\medskip
\textbf{User:}
\texttt{[QUERY]} $\rightarrow$
\texttt{[CRITERIA]} $\rightarrow$
\textbf{\texttt{[EXAMPLES]}} $\rightarrow$
\texttt{[ANSWER]}
\end{minipage}%
}
\caption{Original prompt format. The system prompt announces fixed few-shot
examples, and the user prompt includes an example block containing rationales
and labels.}
\label{fig:original-prompt}
\end{figure}

\begin{figure}[H]
\centering
\fbox{%
\begin{minipage}{0.92\columnwidth}
\small
\textbf{Normalized prompt}

\medskip
\textbf{System:} You are an evaluator of expert-domain scientific writing.
You will receive a query--answer pair, evaluation criteria, and a scoring
rubric. First output your reasoning between \texttt{<reasoning>} and
\texttt{</reasoning>}, followed by the numeric score between
\texttt{<score>} and \texttt{</score>}.

\medskip
\textbf{User:}
\texttt{[QUERY]} $\rightarrow$
\texttt{[CRITERIA]} $\rightarrow$
\texttt{[ANSWER]}
\end{minipage}%
}
\caption{Normalized prompt format. The fixed few-shot examples and the
corresponding system instruction are removed, while all task-specific content
is retained.}
\label{fig:normalized-prompt}
\end{figure}

\clearpage

\begin{figure*}[p]
\centering
\fbox{%
\begin{minipage}{0.95\textwidth}
\raggedright
\footnotesize
\textbf{CoT training example: RevUtil actionability}

\medskip
\textbf{System prompt.}
You are an evaluator of expert-domain scientific writing. You will get a
query-answer pair along with criteria explaining the specific evaluation
aspect and the scoring rubric. You should evaluate whether the answer satisfy
the query based on the given criteria. First output your reasoning enclosed
between \texttt{<reasoning>} and \texttt{</reasoning>}. Then, output your
score enclosed between \texttt{<score>} and \texttt{</score>}. Inside
\texttt{<score>} provide only the numeric score and nothing else.

\medskip
\textbf{User prompt.}

\texttt{[QUERY]:} Your task is to write a review comment for a scientific
paper. The comment should be actionable. Those actions should be clearly
identifiable and concrete.

\medskip
\texttt{[CRITERIA]:} Explicit actions or suggestions are direct or apparent.
Authors can directly identify modifications they should apply to their draft.
Clarification questions should be treated as explicit statements if they give
a direct action. However, implicit actions need to be inferred from the
comment. This includes missing parts that need to be added. Authors can deduce
what needs to be done after reading the comment. For concrete actions, the
authors know exactly what needs to be done and how to apply the action.
However, for vague actions the authors still do not know how to carry out this
action. Scoring rubric is as follows:

1: The comment lacks meaningful information to help authors improve the paper.
Authors do not know what they should do after reading the comment.

2: The comment includes an implicitly stated action or an action that can be
inferred. However, the action itself is vague and lacks detail on how to apply
it.

3: The comment explicitly states an action but is vague on how to execute it.

4: The comment implicitly states an action but concretely states how to
implement the inferred action.

5: The comment contains an explicit action and concrete details on how to
implement it. Authors know exactly how to apply it.

\medskip
\texttt{[ANSWER]:} - line 114: Why not mention that the operator has to be
associative and commutative?

\medskip
\textbf{Gold label.} 5

\medskip
\textbf{Completion.}
\texttt{<reasoning>} The comment explicitly suggests adding a mention that the
operator must be associative and commutative, specifically at line 114. This
is a clear, direct action with a concrete location. The authors know exactly
what to write and where, making it fully actionable and concrete.
\texttt{</reasoning>} \texttt{<score>5</score>}
\end{minipage}%
}
\caption{Complete CoT training example for RevUtil actionability.}
\label{fig:cot-example}
\end{figure*}

\clearpage

\begin{figure*}[p]
\centering
\fbox{%
\begin{minipage}{0.95\textwidth}
\raggedright
\footnotesize
\textbf{Label-only training example: RevUtil actionability}

\medskip
\textbf{System prompt.}
You are an evaluator of expert-domain scientific writing. You will get a
query-answer pair along with criteria explaining the specific evaluation
aspect and the scoring rubric. You should evaluate whether the answer satisfy
the query based on the given criteria. Output your score enclosed between
\texttt{<score>} and \texttt{</score>}. Inside \texttt{<score>} provide only
the numeric score and nothing else. Do not output any reasoning or other text.

\medskip
\textbf{User prompt.}

\texttt{[QUERY]:} Your task is to write a review comment for a scientific
paper. The comment should be actionable. Those actions should be clearly
identifiable and concrete.

\medskip
\texttt{[CRITERIA]:} Explicit actions or suggestions are direct or apparent.
Authors can directly identify modifications they should apply to their draft.
Clarification questions should be treated as explicit statements if they give
a direct action. However, implicit actions need to be inferred from the
comment. This includes missing parts that need to be added. Authors can deduce
what needs to be done after reading the comment. For concrete actions, the
authors know exactly what needs to be done and how to apply the action.
However, for vague actions the authors still do not know how to carry out this
action. Scoring rubric is as follows:

1: The comment lacks meaningful information to help authors improve the paper.
Authors do not know what they should do after reading the comment.

2: The comment includes an implicitly stated action or an action that can be
inferred. However, the action itself is vague and lacks detail on how to apply
it.

3: The comment explicitly states an action but is vague on how to execute it.

4: The comment implicitly states an action but concretely states how to
implement the inferred action.

5: The comment contains an explicit action and concrete details on how to
implement it. Authors know exactly how to apply it.

\medskip
\texttt{[ANSWER]:} - line 114: Why not mention that the operator has to be
associative and commutative?

\medskip
\textbf{Gold label.} 5

\medskip
\textbf{Completion.} \texttt{<score>5</score>}
\end{minipage}%
}
\caption{Complete label-only training example for the same RevUtil
actionability instance as Figure~\ref{fig:cot-example}.}
\label{fig:label-only-example}
\end{figure*}
